% GENERATED FILE. Edit canonical lessons, then run npm run generate:books.
% canonical-source-hash: fnv1a64:e26fff7c7927757d
% canonical-lessons: GU-C513-udyog, GU-C513-kampni, GU-C513-masin, GU-C513-pittal, GU-C513-bhangar

\chapter{On the Factory Floor}
\label{ch:gu-on-the-factory-floor}
\hlchaptermodality{513}

\begin{chapteropening}
\textbf{By the end of this chapter:} \emph{I can talk about industries, companies and machines.}

Complete the last lesson of chapter 513: I can talk about industries, companies and machines.
\end{chapteropening}

\section[\texorpdfstring{\gu{ઉદ્યોગ} (udyog)}{ઉદ્યોગ (udyog)}]{\gu{ઉદ્યોગ} (udyog) — an industry}

\label{lesson:GU-C513-udyog}



\begin{quote}
Before the new one: say the Gujarati for a career, then the Gujarati for a promotion.
\end{quote}

\subsection*{You'll want to know: \gu{ઉદ્યોગ}}
\textbf{\gu{ઉદ્યોગ}} — \emph{udyog} — \textquotedblleft{}an industry\textquotedblright{}.

\begin{quote}
\gu{અમારા} \gu{શહેરમાં} \gu{ઘણા} \gu{મોટા} \gu{ઉદ્યોગ} \gu{છે}.
\end{quote}

\begin{quote}
\emph{amārā śahermā̃ ghaṇā moṭā udyog che.}
\end{quote}

\begin{quote}
There are many big industries in our city.
\end{quote}

A masculine noun.

\begin{cousinweb}
\textbf{\gu{ઉદ્યોગ}} (\emph{udyog}) is a tatsama, taken from Sanskrit \emph{udyoga}, \textquotedblleft{}exertion, an undertaking\textquotedblright{} (Monier-Williams); modern Gujarati uses it for an industry.
\end{cousinweb}

\begin{grammarlens}[title={What you've built}]
One more word for describing what you lived and felt.
\end{grammarlens}

\subsection*{Your turn}
\begin{itemize}
\raggedright
  \item \emph{Say it:} \emph{udyog}
  \item \emph{Say it:} \emph{udyog}, once more
  \item \emph{Say it:} \emph{amārā śahermā̃ ghaṇā moṭā udyog che.} — There are many big industries in our city.
  \item \emph{Recall:} say the Gujarati for a career, then the Gujarati for a promotion, then say \emph{udyog} again
\end{itemize}

\begin{tcolorbox}[breakable,colback=teal!4,colframe=teal!35!black,title={Before you move on}]
What does \gu{ઉદ્યોગ} mean? (\textquotedblleft{}An industry\textquotedblright{}.) Say it once more.
\end{tcolorbox}
\section[\texorpdfstring{\gu{કંપની} (kampnī)}{કંપની (kampnī)}]{\gu{કંપની} (kampnī) — a company (a business)}

\label{lesson:GU-C513-kampni}



\begin{quote}
Before the new one: say the Gujarati for a promotion, then the Gujarati for an industry.
\end{quote}

\subsection*{You'll want to know: \gu{કંપની}}
\textbf{\gu{કંપની}} — \emph{kampnī} — \textquotedblleft{}a company (a business)\textquotedblright{}.

\begin{quote}
\gu{મારા} \gu{ભાઈએ} \gu{એક} \gu{નાની} \gu{કંપની} \gu{શરૂ} \gu{કરી}.
\end{quote}

\begin{quote}
\emph{mārā bhāīe ek nānī kampnī śarū karī.}
\end{quote}

\begin{quote}
My brother started a small company.
\end{quote}

A feminine noun.

\begin{cousinweb}
\textbf{\gu{કંપની}} (\emph{kampnī}) is borrowed from English \emph{company}.
\end{cousinweb}

\begin{grammarlens}[title={What you've built}]
One more word for describing what you lived and felt.
\end{grammarlens}

\subsection*{Your turn}
\begin{itemize}
\raggedright
  \item \emph{Say it:} \emph{kampnī}
  \item \emph{Say it:} \emph{kampnī}, once more
  \item \emph{Say it:} \emph{mārā bhāīe ek nānī kampnī śarū karī.} — My brother started a small company.
  \item \emph{Recall:} say the Gujarati for a promotion, then the Gujarati for an industry, then say \emph{kampnī} again
\end{itemize}

\begin{tcolorbox}[breakable,colback=teal!4,colframe=teal!35!black,title={Before you move on}]
What does \gu{કંપની} mean? (\textquotedblleft{}A company (a business)\textquotedblright{}.) Say it once more.
\end{tcolorbox}
\section[\texorpdfstring{\gu{મશીન} (maśīn)}{મશીન (maśīn)}]{\gu{મશીન} (maśīn) — a machine}

\label{lesson:GU-C513-masin}



\begin{quote}
Before the new one: say the Gujarati for an industry, then the Gujarati for a company.
\end{quote}

\subsection*{You'll want to know: \gu{મશીન}}
\textbf{\gu{મશીન}} — \emph{maśīn} — \textquotedblleft{}a machine\textquotedblright{}.

\begin{quote}
\gu{કારખાનાનું} \gu{એક} \gu{મશીન} \gu{બગડી} \gu{ગયું}.
\end{quote}

\begin{quote}
\emph{kārkhānānũ ek maśīn bagḍī gayũ.}
\end{quote}

\begin{quote}
One of the factory's machines broke down.
\end{quote}

A neuter noun.

\begin{cousinweb}
\textbf{\gu{મશીન}} (\emph{maśīn}) is borrowed from English \emph{machine}.
\end{cousinweb}

\begin{grammarlens}[title={What you've built}]
One more word for describing what you lived and felt.
\end{grammarlens}

\subsection*{Your turn}
\begin{itemize}
\raggedright
  \item \emph{Say it:} \emph{maśīn}
  \item \emph{Say it:} \emph{maśīn}, once more
  \item \emph{Say it:} \emph{kārkhānānũ ek maśīn bagḍī gayũ.} — One of the factory's machines broke down.
  \item \emph{Recall:} say the Gujarati for an industry, then the Gujarati for a company, then say \emph{maśīn} again
\end{itemize}

\begin{tcolorbox}[breakable,colback=teal!4,colframe=teal!35!black,title={Before you move on}]
What does \gu{મશીન} mean? (\textquotedblleft{}A machine\textquotedblright{}.) Say it once more.
\end{tcolorbox}
\section[\texorpdfstring{\gu{પિત્તળ} (pittaḷ)}{પિત્તળ (pittaḷ)}]{\gu{પિત્તળ} (pittaḷ) — brass}

\label{lesson:GU-C513-pittal}



\begin{quote}
Before the new one: say the Gujarati for a company, then the Gujarati for a machine.
\end{quote}

\subsection*{You'll want to know: \gu{પિત્તળ}}
\textbf{\gu{પિત્તળ}} — \emph{pittaḷ} — \textquotedblleft{}brass\textquotedblright{}.

\begin{quote}
\gu{આ} \gu{દીવો} \gu{પિત્તળનો} \gu{છે}.
\end{quote}

\begin{quote}
\emph{ā dīvo pittaḷno che.}
\end{quote}

\begin{quote}
This lamp is made of brass.
\end{quote}

A neuter noun.

\begin{cousinweb}
\textbf{\gu{પિત્તળ}} (\emph{pittaḷ}) is inherited from Sanskrit \emph{pittala}, \textquotedblleft{}brass\textquotedblright{}, and CDIAL lists the Gujarati word (CDIAL 8184).
\end{cousinweb}

\begin{grammarlens}[title={What you've built}]
One more word for describing what you lived and felt.
\end{grammarlens}

\subsection*{Your turn}
\begin{itemize}
\raggedright
  \item \emph{Say it:} \emph{pittaḷ}
  \item \emph{Say it:} \emph{pittaḷ}, once more
  \item \emph{Say it:} \emph{ā dīvo pittaḷno che.} — This lamp is made of brass.
  \item \emph{Recall:} say the Gujarati for a company, then the Gujarati for a machine, then say \emph{pittaḷ} again
\end{itemize}

\begin{tcolorbox}[breakable,colback=teal!4,colframe=teal!35!black,title={Before you move on}]
What does \gu{પિત્તળ} mean? (\textquotedblleft{}Brass\textquotedblright{}.) Say it once more.
\end{tcolorbox}
\section[\texorpdfstring{\gu{ભંગાર} (bhangār)}{ભંગાર (bhangār)}]{\gu{ભંગાર} (bhangār) — scrap}

\label{lesson:GU-C513-bhangar}



\begin{quote}
Before the new one: say the Gujarati for a machine, then the Gujarati for brass.
\end{quote}

\subsection*{You'll want to know: \gu{ભંગાર}}
\textbf{\gu{ભંગાર}} — \emph{bhangār} — \textquotedblleft{}scrap\textquotedblright{}.

\begin{quote}
\gu{આ} \gu{જૂનું} \gu{મશીન} \gu{હવે} \gu{ભંગારમાં} \gu{જશે}.
\end{quote}

\begin{quote}
\emph{ā jūnũ maśīn have bhangārmā̃ jaśe.}
\end{quote}

\begin{quote}
This old machine will now go for scrap.
\end{quote}

A masculine noun.

\begin{cousinweb}
\textbf{\gu{ભંગાર}} (\emph{bhangār}) comes from \emph{*bhaṅgākara}, a form CDIAL reconstructs, and CDIAL lists the Gujarati word as \textquotedblleft{}broken pots\textquotedblright{} (CDIAL 9356); today it is any broken metal or junk sold for scrap.
\end{cousinweb}

\begin{grammarlens}[title={What you've built}]
One more word for describing what you lived and felt.
\end{grammarlens}

\subsection*{Your turn}
\begin{itemize}
\raggedright
  \item \emph{Say it:} \emph{bhangār}
  \item \emph{Say it:} \emph{bhangār}, once more
  \item \emph{Say it:} \emph{ā jūnũ maśīn have bhangārmā̃ jaśe.} — This old machine will now go for scrap.
  \item \emph{Recall:} say the Gujarati for a machine, then the Gujarati for brass, then say \emph{bhangār} again
\end{itemize}

\begin{tcolorbox}[breakable,colback=teal!4,colframe=teal!35!black,title={Before you move on}]
What does \gu{ભંગાર} mean? (\textquotedblleft{}Scrap\textquotedblright{}.) Say it once more.
\end{tcolorbox}
